\scanpage{15}
% Pencil foliation: 13, upper right; 14, lower left.
12) Structures supplémentaires sur le groupe de Galois
motivique : \inserted{structures associées aux}
graduations, filtrations, au motif de Tate.

On choisit une sous-\inserted{$\otimes$-}catégorie graduée
\deleted{la catégorie}
$C$ de $\mathcal{M}(S)$, avec un foncteur fibre $F$,
d'où groupe de Galois motivique $G$
\inserted{(à coeff dans $K$)}.

a) Une graduation de type $\Pi$ du foncteur fibre,
compatible avec $\otimes$, est équivalente à la donnée
d'un hom
\diagram{\begin{tikzcd}
T=D_K(\unsure{\Pi|\pi})\arrow[r]&G
\end{tikzcd}}
La graduation provient d'une graduation fonctorielle
de $M\inserted{\in C}$ lui-même ssi l'hom $T\longrightarrow G$ est
central. Par suite :

1°) La graduation \inserted{canonique} de type $\mathbb{Z}$ des objets de
$C$ définit un hom canonique
\diagram{\begin{tikzcd}
\unsure{\gamma|g}:\mathbb{G}_{\mathrm m,K}\arrow[r]&G
\end{tikzcd}}

2°) Si $F=T_{\mathrm{Hdg}}$, la bigraduation de ce
foncteur correspond à un hom
\deleted{\diagram{\begin{tikzcd}
D_K(\mathbb{Z}^2)=\mathbb{G}_{\mathrm m}^2\arrow[r]&G
\end{tikzcd}}
donc à deux \unsure{h}}
\diagram{\begin{tikzcd}
i:\mathbb{G}_{\mathrm m,K}^2\arrow[r]&G
\end{tikzcd}}
donc à deux hom commutant l'un à
l'autre
\diagram{\begin{tikzcd}
i_1,i_2:\mathbb{G}_{\mathrm m}\arrow[r]&G
\end{tikzcd}}
Comme la graduation totale associée est celle
de 1°), on trouve
\[
 i_1(\lambda)i_2(\lambda)=\unsure{\gamma|g}(\lambda)
\]
(ce qui entraîne déjà que $i_1,i_2$ commutent \dots).
